\documentclass[11pt,a4paper]{article}

% -------------------------------------------------
% Page layout and typography
% -------------------------------------------------
\usepackage[
    top=22mm,
    bottom=24mm,
    left=25mm,
    right=25mm
]{geometry}

\usepackage[T1]{fontenc}
\usepackage{lmodern}
\usepackage{microtype}
\usepackage{parskip}
\usepackage{ragged2e}

% -------------------------------------------------
% Tables and lists
% -------------------------------------------------
\usepackage{booktabs}
\usepackage{siunitx}
\usepackage{enumitem}

\sisetup{
    group-separator = {,},
    group-minimum-digits = 4,
    table-number-alignment = right,
    round-mode = places,
    round-precision = 2
}

% -------------------------------------------------
% Headings, colours, and page style
% -------------------------------------------------
\usepackage{xcolor}
\usepackage{titlesec}
\usepackage{fancyhdr}

\definecolor{headingblue}{HTML}{17365D}
\definecolor{linkblue}{HTML}{1F5A94}
\definecolor{lightgray}{HTML}{666666}

\titleformat{\section}
  {\Large\bfseries\color{headingblue}}
  {\thesection}
  {0.7em}
  {}

\titlespacing*{\section}
  {0pt}
  {2.2ex plus 0.5ex minus 0.2ex}
  {0.8ex}

\pagestyle{fancy}
\fancyhf{}
\fancyhead[L]{\small\color{lightgray}COMP3331/9331}
\fancyhead[R]{\small\color{lightgray}DNS Resolver Assignment Report}
\fancyfoot[C]{\thepage}

\setlength{\headheight}{14pt}
\renewcommand{\headrulewidth}{0.4pt}
\renewcommand{\footrulewidth}{0pt}

% -------------------------------------------------
% Links and code formatting
% -------------------------------------------------
\usepackage{xurl}
\usepackage[
    colorlinks=true,
    linkcolor=linkblue,
    urlcolor=linkblue,
    citecolor=linkblue,
    pdfauthor={},
    pdftitle={COMP3331/9331 DNS Resolver Assignment Report}
]{hyperref}

\urlstyle{same}

\newcommand{\code}[1]{\texttt{#1}}
\newcommand{\program}[1]{\texttt{\textbf{#1}}}

% -------------------------------------------------
% Document metadata
% -------------------------------------------------
\title{
    \vspace{-1.5cm}
    {\Huge\bfseries\color{headingblue}
    DNS Resolver Assignment Report}
    \\[0.5em]
    {\Large COMP3331/9331}
}
\author{Owais Wangde\\z5690720}
\date{\small 2026 Term 2}

\begin{document}

\maketitle
\vspace{-1.5em}

\begin{center}
    \rule{0.82\textwidth}{0.5pt}
\end{center}

\vspace{1em}

% =================================================
\section{Code organisation}

The submission consists of two main programs.

\program{parser.py} reads a single DNS message from a binary file and displays
the decoded message in a readable format. It operates entirely on local files
and does not perform any network communication.

Its program flow is \code{main()} $\rightarrow$ read message file
$\rightarrow$ \code{print\_message()} $\rightarrow$ \code{parse\_message()}.
The parser then prints the header, Questions, Answers, Authority, and Additional
sections in wire order.

\program{resolver.py} implements the UDP DNS resolver. At startup, it reads the
root hints file, binds to \code{127.0.0.1}, and handles each client datagram in
a separate thread. Its flow is \code{main()} $\rightarrow$ \code{serve()}
$\rightarrow$ \code{Resolver.response()}. For each query, \code{response()}
checks the root-hints-local cases, searches the cache, and calls
\code{resolve()} when no usable cached result is available. It then calls
\code{build\_client\_response()} to send the reply.

The programs share code through the \program{parser} module. \program{resolver.py}
imports \code{parser.py} and calls \code{parser.parse\_message()} to decode both
client queries and upstream DNS responses. It also uses the shared
\code{Question}, \code{ResourceRecord}, and \code{DNSParseError} definitions.
This keeps the parser's bounds checks, compression handling, and RDLENGTH rules
the same in both programs.

The automated tests are grouped by parser, resolver-core, black-box, and
integration behaviour. The integration suite includes a deterministic UDP
concurrency stress test that exercises the resolver and its shared cache without
depending on public DNS.

% =================================================
\section{Data structures and algorithms}

The \code{Question} structure stores a domain name, query type, and query class.
The resource-record structure \code{ResourceRecord} stores:

\begin{itemize}[leftmargin=2em, itemsep=0.2em]
    \item the owner name;
    \item record type and class;
    \item time to live (TTL);
    \item RDLENGTH;
    \item decoded record data; and
    \item the original RDATA bytes.
\end{itemize}

The resolver also uses \code{DNSRecord}, a smaller record structure for the
supported DNS record types.

The DNS message representation is the header values, a list of \code{Question}
objects, and three ordered lists of \code{ResourceRecord} objects for the Answer,
Authority, and Additional sections. \code{DNSMessage} keeps the input bytes;
the parsing functions pass an offset forward as each field is decoded. Names are
represented as strings with a trailing dot. The encoder maps case-normalised
name suffixes to packet offsets, allowing a mixed-case owner to point to an
earlier name while retaining the original Question spelling.

The parser reads the DNS header, question section, and all resource records in
the Answer, Authority, and Additional sections. It supports the following
record types:

\begin{center}
    \code{A}, \code{NS}, \code{CNAME}, \code{MX}, and \code{PTR}.
\end{center}

The parser follows DNS compression pointers and detects pointer loops. A domain
name is limited to 20 pointer jumps, and the parser verifies the end of every
field before reading it. The RDLENGTH value determines the endpoint of each
RDATA field. If an RDATA field is malformed, the parser reports the error while
preserving the correct position for reading the following record.

The resolver begins with the root name servers and their matching A records
from the root hints file. Candidate servers are tried in order. A referral uses
IN-class NS records in the Authority section and matching IN-class A glue in
the Additional section. When usable glue is unavailable, the resolver performs
a separate lookup from the root to obtain the referred name server's A record.

Referral state is kept in a \code{Budget} object for one client request. It
contains the deadline and counters for upstream attempts, referral levels, and
CNAME records. The current candidate server list and CNAME chain are local
variables in the resolution function. They are not shared between clients.

Cache entries are indexed using the following key:

\begin{center}
    \code{(lowercase fully qualified name, type, class)}
\end{center}

Each cache entry stores an expiry time calculated using
\code{time.monotonic()}. The TTL returned to a client decreases as time passes.
Records with a TTL of zero are not cached. For an alias lookup, the cache can
return a complete CNAME chain followed by the requested final record set.

The cache is the only shared mutable resolver state. A \code{threading.Lock}
protects cache reads and updates. Each UDP client request has its own worker
thread, query bytes, client address and port, transaction ID, \code{Budget},
candidate list, and CNAME chain. This prevents one lookup from changing another
lookup's resolution state.

% =================================================
\section{Design decisions and limitations}

\textbf{Parsing and matching.} All client and upstream packets use the same
bounds-checked parser. It implements the DNS header, section layout, RDLENGTH
boundaries, and compression pointers described by RFC 1035 [2, Section~4.1].
A malformed client packet receives no reply. For upstream UDP, datagrams with a
wrong source endpoint, ID, opcode, or Question are unrelated and ignored, as
motivated by RFC 1035 [2, Section~7.3]. Once those identity fields match, TC,
malformed structure or supported RDATA, an unsupported RCODE, or a socket error
fails that candidate immediately; resolution continues with the next server
without waiting out the remaining timeout.

\textbf{Iterative resolution.} A cache miss starts from ordered root-hint
addresses and sends non-recursive IPv4 UDP queries with fresh random 16-bit IDs.
The answer, name-error, closer-delegation, CNAME, and server-failure states follow
RFC 1034 [1, Section~5.3.3]. A referral is accepted only when an IN-class
Authority NS RRset moves to a longer relevant suffix; only matching IN-class A
glue is used. Otherwise, a nested A lookup resolves the nameserver.

\textbf{Aliases and negative answers.} For A, NS, MX, and PTR, CNAME RRsets are
kept in wire order before the final type. Targets must be consistent; loops and
more than ten retained CNAME records cause SERVFAIL. A direct CNAME query returns
the RRset without chasing it, consistent with RFC 1034 [1, Sections~4.3.2
and~5.2.2]. Authoritative NXDOMAIN is final, but reachable CNAME links in that
same packet are retained along with links from earlier packets. AA is copied only
when the complete positive, NODATA, or NXDOMAIN result came from one authoritative
packet; cached and cross-packet results clear it.

\textbf{Caching and response construction.} Positive RRsets use lowercase
\code{(name, type, class)} keys, absolute monotonic expiry, decreasing returned
TTLs, and no TTL-zero reuse, following RFC 1035 [2, Section~6.1.3]. The response
encoder preserves the client's ID, opcode, RD, and Question spelling, sets QR and
RA, derives exact counts, and safely compresses case-insensitive suffix matches.
At the classic UDP limit in RFC 1035 [2, Section~4.2.1], this assignment requires
TC=0, so an encoding over 512 bytes becomes an empty SERVFAIL.

\textbf{Limits and concurrency.} Each request has a
\code{min(30, 50 * timeout)} deadline and caps of 50 upstream attempts, ten
referrals, and ten retained CNAME records. Independent worker threads prevent a
slow lookup from serialising other UDP clients, targeting RFC 1035 [2,
Section~6.1.1], while a lock protects the only shared mutable object, the cache.
The implementation uses only Python's standard library. It does not implement
TCP fallback, IPv6, EDNS(0), DNSSEC validation, negative or persistent caching,
or record types outside the assignment. A production resolver would also bound
worker threads and request rates.

% =================================================
\section{Performance report}

On 30 July 2026, a raw-UDP client on the same MacBook Air tested the submitted
resolver and Cloudflare \code{1.1.1.1}. Every exchange used a fresh transaction
ID, the fixed order below, and a three-second client timeout. Latency ran from
\code{sendto} to the matching \code{recvfrom}; total is batch wall-clock time.

\begin{description}[leftmargin=5.5em, style=nextline, itemsep=0.1em]
    \item[A] \code{example.com}, \code{www.iana.org},
    \code{www.python.org}, \code{www.unsw.edu.au}
    \item[NS] \code{example.com}, \code{iana.org}, \code{python.org},
    \code{unsw.edu.au}
    \item[MX] \code{example.com}, \code{gmail.com}, \code{outlook.com},
    \code{unsw.edu.au}
    \item[CNAME] \code{www.python.org}, \code{www.wikipedia.org},
    \code{www.microsoft.com}, \code{www.github.com}
    \item[PTR] reversed names for \code{8.8.8.8}, \code{1.1.1.1},
    \code{4.2.2.4}, and \code{8.8.4.4}
\end{description}

One resolver process handled cold then warm sequential batches; a fresh process
handled cold then warm concurrent batches. The public batches used the same
client and order, with RD set for Cloudflare. Thus only the local resolver's
cache state was controlled; Cloudflare's pre-existing cache state is unknown.

\begin{table}[htbp]
    \centering
    \scriptsize
    \setlength{\tabcolsep}{2.2pt}
    \caption{Fixed 20-query raw-UDP comparison (times in ms)}
    \label{tab:performance}
    \vspace{0.3em}

    \begin{tabular}{
        lrrrr
        S[table-format=5.2]
        S[table-format=4.2]
        S[table-format=4.2]
        S[table-format=4.2]
    }
        \toprule
        \textbf{Batch}
        & \textbf{Q}
        & \textbf{OK}
        & \textbf{RC}
        & \textbf{TO}
        & {\textbf{Total}}
        & {\textbf{Mean}}
        & {\textbf{Median}}
        & {\textbf{p95}} \\
        \midrule
        Local cold seq. & 20 & 18 & 2 & 0
            & 16865.74 & 843.26 & 301.81 & 2752.44 \\
        Local warm seq. & 20 & 18 & 2 & 0
            & 5924.87 & 296.23 & 0.12 & 2939.55 \\
        Local cold conc. & 20 & 18 & 2 & 0
            & 2847.37 & 825.48 & 334.41 & 2683.31 \\
        Local warm conc. & 20 & 18 & 2 & 0
            & 3132.49 & 358.20 & 69.69 & 2855.03 \\
        Cloudflare seq. & 20 & 20 & 0 & 0
            & 927.34 & 46.35 & 11.62 & 168.35 \\
        Cloudflare conc. & 20 & 20 & 0 & 0
            & 207.32 & 55.86 & 40.13 & 178.41 \\
        \bottomrule
    \end{tabular}
\end{table}

All 120 exchanges received replies. The two local non-zero results in every
batch were SERVFAIL for \code{www.iana.org A} and
\code{1.1.1.1.in-addr.arpa PTR}; they were not timeouts. Warm sequential median
fell from 301.81 ms to 0.12 ms because many positive answers were cached.
Concurrency reduced cold batch wall time from 16.87 s to 2.85 s. Warm concurrent
median also improved, although its total was slightly above cold concurrent due
to Internet variation and uncached/NODATA work. Cloudflare was faster, but it is
a distributed recursive service with unknown cache state, so this is contextual
rather than a controlled implementation-to-implementation comparison.

\subsection*{Concurrency stress test}

On 30 July 2026, the concurrency stress test in
\path{tests/test_resolver_integration.py} was run on the same host. It uses a
local scripted authoritative UDP server, so the result is repeatable and
independent of Internet conditions. The test sends five waves of 40 clients
released simultaneously: one cold-cache wave and four warm-cache waves, for 200
requests in total. Each request has a unique transaction ID, while the 40 names
and expected addresses repeat across waves. The upstream server adds
deterministic delays of 0--21 ms to reorder cold responses.

\begin{table}[htbp]
    \centering
    \caption{Concurrent resolver stress-test result}
    \label{tab:concurrency-stress}
    \vspace{0.5em}

    \begin{tabular}{lr}
        \toprule
        \textbf{Measure} & \textbf{Result} \\
        \midrule
        Client requests attempted & 200 \\
        Responses received and validated & 200 \\
        Timeouts or client errors & 0 \\
        Transaction-ID or answer mismatches & 0 \\
        Upstream queries & 40 \\
        \bottomrule
    \end{tabular}
\end{table}

The 40 upstream queries came from the cold wave; all 160 later requests were
served from the shared cache. The assertions verify every response ID and A
record rather than only counting replies. This stresses simultaneous socket
handling, per-request state isolation, concurrent cache writes, and concurrent
cache reads. The existing slow-versus-fast test separately confirms that a
300 ms upstream delay does not serialise an unrelated request.

% =================================================
\section{Security discussion}

\textbf{DNS cache poisoning.} An attacker may send a forged UDP response before
the legitimate upstream response arrives.

For each upstream request, this resolver generates a new random transaction ID.
It also validates the response's source IP address, source port, transaction ID,
flags, and Question section. These checks reduce the likelihood that a forged
response will be accepted. However, they do not provide cryptographic
authentication. DNSSEC validation is outside the scope of the assignment.

\textbf{Resource exhaustion and amplification.} A slow upstream server can
occupy a worker thread, while many client packets can create threads. A resolver
that is reachable from the Internet can also be used to reflect larger DNS
responses towards another host.

The resolver limits the resources consumed by a single request through
timeouts and limits on upstream attempts, referrals, and CNAME records. It also
rejects malformed domain names and malformed RDATA fields safely.

Finally, the resolver binds only to \code{127.0.0.1}. This prevents remote
Internet clients from using it as an open recursive resolver or as part of a DNS
amplification attack.

% =================================================
\section{References}

\begin{enumerate}[leftmargin=2em, label={[\arabic*]}, itemsep=0.8em]
    \item P. Mockapetris,
    \textit{Domain Names: Concepts and Facilities},
    RFC 1034, November 1987.\\
    \url{https://datatracker.ietf.org/doc/html/rfc1034}

    \item P. Mockapetris,
    \textit{Domain Names: Implementation and Specification},
    RFC 1035, November 1987.\\
    \url{https://datatracker.ietf.org/doc/html/rfc1035}

    \item Internet Assigned Numbers Authority,
    \textit{Domain Name System (DNS) Parameters}.\\
    \url{https://www.iana.org/assignments/dns-parameters/dns-parameters.xhtml}

    \item COMP3331/9331,
    \textit{Assignment: Building a DNS Resolver},
    2026 Term 2, version v1.

    \item R. Elz and R. Bush,
    \textit{Clarifications to the DNS Specification},
    RFC 2181, July 1997.\\
    \url{https://datatracker.ietf.org/doc/html/rfc2181}

    \item D. Andrews,
    \textit{Negative Caching of DNS Queries (DNS NCACHE)},
    RFC 2308, March 1998.\\
    \url{https://datatracker.ietf.org/doc/html/rfc2308}

    \item R. Arends et al.,
    \textit{DNS Security Introduction and Requirements},
    RFC 4033, March 2005.\\
    \url{https://datatracker.ietf.org/doc/html/rfc4033}
\end{enumerate}

% =================================================
\section{AI interaction summary}

Codex was used between 15 and 30 July 2026 as a learning, design-review, and
testing aid. During 15--17 July, I used it to clarify DNS name compression,
including pointer offsets, nested pointers, loop and bounds checks, RDLENGTH
boundaries, and retaining the outer-field offset after following a pointer. On
16 July, AI produced a sample name-decoding algorithm, which I rewrote and
adapted to the \program{parser.py} structure rather than copying directly.

From 19--26 July, I discussed how to organise \program{resolver.py} and reasoned
through later assignment stages. Stage 2 discussion covered root-hints parsing,
local root NS/A responses, flags, counts, and failure handling; the 22 July
interaction was conceptual and generated no submitted implementation. Stage 3
discussion covered iterative referrals, IN-class glue, nested nameserver-address
lookups, CNAME and negative answers, caching, timeouts, response size, and
concurrency.

On 27 July, AI suggested parser, resolver, malformed-packet, cache, failover, and
concurrency tests. On 30 July, it helped design the deterministic 200-request
stress test, identify relevant RFC 1034/1035 sections after tutor feedback, and
audit edge cases. This affected \program{tests/}, \program{report.tex}, and
targeted parts of \program{resolver.py}.

I retained ideas that matched the specification, rewrote or simplified proposed
algorithms, and rejected external DNS libraries, changed assessed interfaces,
public recursive resolution, TCP fallback, and other out-of-scope features. I
checked retained work against the assignment, RFCs, released packets, automated
and live tests, repeated stress runs, and PDF inspection. I reviewed and can
explain the parser, compression, iterative resolver, cache, and concurrency
components in the viva. Detailed dated interactions are in
\program{ai-log.txt}.

\end{document}
